\documentclass{ximera}
\typeout{Start loading xmPreamble.tex}%

% Add here extra macro's that are loaded automatically by all documents of claas 'ximera' or 'xourse' in this repo

%%
%%  Example:
%%
% \newcommand{\R}{\mathbb{R}}
\usepackage{tikz}
\usetikzlibrary{positioning, arrows.meta}
\usepackage{multicol}
\usepackage{enumitem}
\usepackage{caption}
\usepackage{amsmath}
\usepackage{amsthm}

\title{Continuous-Time Markov Processes}
\begin{document}
\begin{abstract}
An introduction to continuous-time Markov processes, including the generator matrix, holding times, and the relationship between the generator and the dynamics of the process.
\end{abstract}
\maketitle

\section*{Continuous-Time Markov Processes}

Last time we proved that $P(s+t) = P(s)P(t)$, also called the Chapman-Kolmogorov equation. For semigroups of real-valued functions $\mathbb{R} \to \mathbb{R}$, one can show (exercise) that
\[
f(s+t) = f(s)f(t) \implies f(t) = e^{\alpha t}
\]
where $\alpha = f'(0)$. Something similar holds for semigroups of stochastic matrices:
\[
P(t) = \exp(Qt), \quad \text{where } Q = P'(0),
\]
and the matrix exponential is defined as
\[
\exp(A) = I + A + \frac{1}{2}A^2 + \cdots
\]
We will see that the matrix $Q$ describes $X_t$ in words.

\begin{theorem}
For any $x \in \mathcal{X}$, the function $P_{xx}(t)$ is differentiable at $t = 0$ (though the derivative might be infinite). Furthermore,
\[
\lim_{t \to 0^+} \frac{P_{xx}(t) - 1}{t}
\]
exists. (Note: $P_{xx}(0) = 1$ since $P(0) = \text{Id}$.)
\end{theorem}

\begin{proof}[Proof (Outline)]
Using that $P(t)$ is stochastic and satisfies the semigroup property,
\[
1 \geq P_{xx}(s+t) = \sum_{y \in \mathcal{X}} P_{xy}(s) P_{yx}(t) \geq P_{xx}(s) P_{xx}(t).
\]
Set $\phi(t) = -\log P_{xx}(t)$, which is subadditive: $\phi(s+t) \leq \phi(s) + \phi(t)$. Define
\[
q = \sup_{t > 0} \frac{\phi(t) - \phi(0)}{t} \in [0, \infty].
\]
Note $\phi(0) = 0$ and $q \geq 0$ by stochasticity. We need to show $\liminf_{t \to 0^+} \frac{\phi(t)}{t} \geq q$. For any $q' < q$, there exists $s > 0$ such that $\frac{\phi(s)}{s} > q'$. Writing $s = nt + h$ with $n \in \mathbb{N}$ and $0 < h < t$, and using $\phi(s) \leq n\phi(t) + \phi(h)$ completes the argument.
\end{proof}

\begin{definition}
If $P'_{xx}(0) = 0$, we call $x$ an \textbf{absorbing state}. If $P'_{xx}(0) = \infty$, we call $x$ an \textbf{instantaneous state}. An absorbing state is one the process never leaves. An instantaneous state is one the process leaves immediately.
\end{definition}

In practice we assume there are no instantaneous states. This assumption has an important consequence:

\begin{proposition}
If $P'_{xx}(0) \neq \infty$ for all $x \in \mathcal{X}$, then $P'_{xy}(0)$ exists and is finite for all $x, y \in \mathcal{X}$.
\end{proposition}

\section*{Generator of a Markov Process}

\begin{definition}
The \textbf{generator} of the Markov process is $Q = P'(0)$.
\end{definition}

Since $P(t)$ has rows summing to 1, $Q = \lim_{t \to 0^+} \frac{P(t) - \text{Id}}{t}$ has rows summing to 0. For $x \neq y$, we write $q_{xy} = Q_{xy}$, and define $q_x = -P'_{xx}(0)$.

\begin{definition}
The \textbf{holding time} at state $x$ is the amount of time spent at $x$ before leaving. Rigorously,
\[
T_x = \inf\{t \geq 0 : X_t \neq x\},
\]
assuming $X_0 = x$.
\end{definition}

\begin{proposition}
\begin{enumerate}[label=(\arabic*)]
\item $P(T_x > t \mid X_0 = x) = e^{-q_x t}$, so $T_x \sim \text{Exp}(q_x)$: the holding time is exponentially distributed with parameter $q_x$.
\item If $q_x \neq 0$, then for $y \neq x$,
\[
P(X_{T_x} = y \mid X_0 = x) = \frac{q_{xy}}{q_x}.
\]
That is, when the process first leaves $x$, it goes to $y$ with probability $q_{xy}/q_x$.
\end{enumerate}
\end{proposition}

\begin{proof}
(1) We use the notation $P_x(\cdot) = P(\cdot \mid X_0 = x)$. First we show the memoryless property:
\[
P_x(T_x > s + t) = P_x(T_x > s) P_x(T_x > t).
\]
The left-hand side equals
\[
P_x(X_u = x\ \forall u \in [0, s+t]).
\]
Conditioning on $\{X_u = x\ \forall u \in [0,t]\}$ and applying the Markov property gives
\[
P_x(T_x > s) P_x(T_x > t),
\]
which is the right-hand side.

Next, we show $P(T_x > t) = e^{-q_x t}$ using the skeleton of the Markov process. By continuity of probability,
\[
P(T_x > t) = P(X_u = x\ \forall u \in [0,t]) = \lim_{n \to \infty} P_{xx}\!\left(\tfrac{t}{n}\right)^n.
\]
Then
\[
-\frac{1}{t} \log P(T_x > t) = -\lim_{\tau \to 0} \frac{\log P_{xx}(\tau)}{\tau} = -\frac{d}{d\tau}\log P_{xx}(\tau)\Big|_{\tau=0} = \frac{-P'_{xx}(0)}{P_{xx}(0)} = q_x,
\]
which proves (1).

(2) Define
\[
R_{xy}(h) = P(X_{t+h} = y \mid X_t = x,\ X_{t+h} \neq x).
\]
As $h \to 0$, this converges to $P(X_{T_x} = y \mid X_0 = x)$. Using the identity $P(A \mid B \cap C) = \frac{P(A \cap B \mid C)}{P(B \mid C)}$,
\[
R_{xy}(h) = \frac{P_{xy}(h)}{1 - P_{xx}(h)}.
\]
Dividing numerator and denominator by $h$ and taking $h \to 0$:
\[
\lim_{h \to 0} R_{xy}(h) = \frac{P'_{xy}(0)}{-P'_{xx}(0)} = \frac{q_{xy}}{q_x}.
\]
\qedhere
\end{proof}

\end{document}